\begin{proofbox}

\textbf{\textcolor{CProof}{证明前统一约定}}

设 $F_1(-c,0)$、$F_2(c,0)$，
$P(x_0,y_0)$ 在双曲线
$\frac{x^2}{a^2}-\frac{y^2}{b^2}=1$ 上，
$c^2=a^2+b^2$，$e=\frac ca$。
记 $d_1=|PF_1|$、$d_2=|PF_2|$。
已知双曲线定义
\[
 |d_1-d_2|=2a.
\]
又由两点间距离公式（注意这里\textbf{只相减，不展开所有项}）得到
\[
 \boxed{d_1^2-d_2^2
 =(x_0+c)^2-(x_0-c)^2=4cx_0.}\tag{*}
\]

\medskip
\textbf{\textcolor{CProof}{方法一：最快证明（定义 $+$ 平方差，首推）}}

\textbf{核心：} 已知两距离的差，再用平方差求出两距离的和；
一加一减即可。\textbf{无需代入双曲线方程消元。}

\begin{description}[style=nextline,leftmargin=2.8em,labelsep=0.5em,itemsep=0.45em]
 \item[\textbf{右支 $x_0\ge a$：}]
  此时到左焦点更远，$d_1-d_2=2a$。
  由 $(*)$：
  \[
   (d_1-d_2)(d_1+d_2)=4cx_0
   \ \Longrightarrow\ d_1+d_2=\frac{4cx_0}{2a}=2ex_0.
  \]
  因而
  \[
   \boxed{d_1=\frac{2ex_0+2a}{2}=ex_0+a},\qquad
   \boxed{d_2=\frac{2ex_0-2a}{2}=ex_0-a}.
  \]
 \item[\textbf{左支 $x_0\le-a$：}]
  此时到右焦点更远，$d_2-d_1=2a$，即 $d_1-d_2=-2a$。
  同理由 $(*)$：
  \[
   -2a(d_1+d_2)=4cx_0
   \ \Longrightarrow\ d_1+d_2=-2ex_0.
  \]
  故
  \[
   \boxed{d_1=-ex_0-a},\qquad
   \boxed{d_2=-ex_0+a}.
  \]
\end{description}
\textbf{方法优势：} 只用双曲线定义、两点距离的平方差、加减消元；
\textbf{速度最快，考场优先采用。}

\medskip
\textbf{\textcolor{CProof}{方法二：最低门槛证明（只设 $t$ 一个未知量）}}

\textbf{为什么这样设？} 两段长度相差 $2a$；
与其分别设两段长，不如先设它们的\textbf{中间值 $t$}。
较长一段为 $t+a$，较短一段为 $t-a$，只有一个未知量。

\begin{description}[style=nextline,leftmargin=2.8em,labelsep=0.5em,itemsep=0.5em]
 \item[\textbf{右支：}]
  $F_1$ 在左边，因而 $PF_1$ 较长。令
  \[
    d_1=t+a,\qquad d_2=t-a\quad(t\ge a).
  \]
  两边各自平方再相减，左边用\textbf{平方差}：
  \[
   \underbrace{(t+a)^2-(t-a)^2}_{4at}
   =\underbrace{\bigl[(x_0+c)^2+y_0^2\bigr]
             -\bigl[(x_0-c)^2+y_0^2\bigr]}_{4cx_0}.
  \]
  所以 $4at=4cx_0$，只剩\textbf{一个}未知数：
  \[
     \boxed{t=\frac ca x_0=ex_0}.
  \]
  代回设法：$d_1=ex_0+a$，$d_2=ex_0-a$。
 \item[\textbf{左支：}]
  $F_2$ 在右边，因而 $PF_2$ 较长。令
  \[
    d_1=t-a,\qquad d_2=t+a\quad(t\ge a).
  \]
  同样两式平方相减：
  \[
    (t-a)^2-(t+a)^2=4cx_0
    \ \Longrightarrow\ -4at=4cx_0,
  \]
  因此 $t=-ex_0=e|x_0|$，进而
  \[
     d_1=-ex_0-a,\qquad d_2=-ex_0+a.
  \]
\end{description}
\textbf{统一理解：} 两距离的中间值始终为
\[
 \boxed{t=\frac{d_1+d_2}{2}=e|x_0|};\qquad
 \boxed{\text{较短}=t-a,\ \text{较长}=t+a}.
\]
\textbf{方法优势：} 完全顺着“相差 $2a$”去设，\textbf{只设 $t$ 一个变量}，
看得见每一步的几何含义；适合第一次学或考后复盘。

\medskip
\textbf{\textcolor{CProof}{方法三：常规坐标证明（标准方程 $+$ 配完全平方）}}

\textbf{核心：} 利用曲线方程消去 $y_0^2$，再开平方。
这是最通用、也最适合检查代数细节的办法。

\begin{description}[style=nextline,leftmargin=2.8em,labelsep=0.5em,itemsep=0.45em]
 \item[\textbf{第一步（消去 $y_0^2$）：}]
   $P$ 在双曲线上，故
   \[
    y_0^2=\frac{b^2}{a^2}(x_0^2-a^2).
   \]
 \item[\textbf{第二步（计算左焦半径的平方）：}]
   \[
   \begin{aligned}
    d_1^2&=(x_0+c)^2+y_0^2\\
    &=x_0^2+2cx_0+c^2+\frac{b^2}{a^2}x_0^2-b^2\\
    &=\frac{a^2+b^2}{a^2}x_0^2+2cx_0+(c^2-b^2)\\
    &=\frac{c^2}{a^2}x_0^2+2cx_0+a^2\\
    &=(ex_0+a)^2.
   \end{aligned}
   \]
   其中使用 $c^2=a^2+b^2$ 及 $ex_0=\frac ca x_0$。
 \item[\textbf{第三步（计算右焦半径）：}]
   只需把 $(x_0+c)^2$ 改为 $(x_0-c)^2$，得到
   \[
    d_2^2=(ex_0-a)^2.
   \]
 \item[\textbf{第四步（开方，绝对值不能丢）：}]
   \[
     \boxed{d_1=\sqrt{(ex_0+a)^2}=|ex_0+a|},\qquad
     \boxed{d_2=\sqrt{(ex_0-a)^2}=|ex_0-a|}.
   \]
\end{description}
\textbf{方法优势：} 不用提前判断哪段更长，一次推出\textbf{带绝对值的统一式}；
代数计算较长，但适用性强。

\medskip
\textbf{\textcolor{CProof}{方法四：第二定义的几何速证（选学）}}

\textbf{前提：} 已经学过双曲线的\textbf{焦点--准线第二定义}：
$|PF|/d(P,\ell)=e$，其中 $\ell$ 是与焦点对应的准线。
本法是对该定义的直接应用，\textbf{不把第二定义当作无需证明的新公理}。

对应 $F_1,F_2$ 的准线分别为
\[
 \ell_1:\ x=-\frac{a^2}{c}=-\frac ae,
 \qquad
 \ell_2:\ x=\frac{a^2}{c}=\frac ae.
\]
对双曲线上任意点 $P(x_0,y_0)$，
\[
 \begin{aligned}
 |PF_1|&=e\,d(P,\ell_1)
   =e\left|x_0+\frac ae\right|=|ex_0+a|,\\
 |PF_2|&=e\,d(P,\ell_2)
   =e\left|x_0-\frac ae\right|=|ex_0-a|.
 \end{aligned}
\]
\textbf{方法优势：} 若熟悉第二定义，这是\textbf{几何意义最直接}的一种证明；
若尚未学第二定义，优先使用方法二。

\medskip
\textbf{\textcolor{CProof}{结果自检（三个不用重算的检查）}}
\begin{itemize}[leftmargin=*,itemsep=0.35em,topsep=0.35em]
 \item \textbf{差：} $\bigl||PF_1|-|PF_2|\bigr|=2a$。
 \item \textbf{和：} $|PF_1|+|PF_2|=2e|x_0|$。
 \item \textbf{顶点：} $x_0=a$ 时 $|PF_1|=c+a$，$|PF_2|=c-a$；
      $x_0=-a$ 时两者互换。
\end{itemize}

\end{proofbox}
